\section{结论}

我们建立了从基础调度到全要素集成的跨区域调度的四层递进的优化模型。从单区域极差最小化出发，逐步引入碳排放约束、储能充放电与新能源基荷预填（定义见 §7.1），最终在半耦合MILP框架中完成多区域算力、储能和电力的协同优化。全文选用MILP、LP与LNS启发式为主要求解工具。最后通过灵敏度分析验证了核心结论的稳健性。

最终我们得出三点关键发现：
\begin{enumerate}
\item 算力调度优化空间高度依赖低价低碳区的电力与算力承载力，把延迟容忍任务从高电价高碳区迁往低价低碳区，只用约2.0\%的成本就能换来约32.7\%的碳排放下降。西部地区电价与碳强度双低，任务迁移一举两得。
\item 储能的作用不是锦上添花，它不仅能峰谷套利，在F区等区域更是保证功率平衡的必需品，评估时不能把储能省下的钱和低价时段买电省下的钱一视同仁，否则会高估储能的价值。
\item 全要素集成模型规模大、难以直接求解，基荷预填以不到0.88\%的精度损失把变量空间压缩约$10^5$倍，让大规模MILP能在数秒内求解。虽然赛题未要求，但以精度换速度在实时调度领域意义重大。
\end{enumerate}

最后，灵敏度分析表明，核心指标对关键参数扰动的方向正确、量级可控。我们的模型在合理参数区间内具有稳健性。其方法论框架可推广至其他涉及多区域资源协同、碳约束与新能源消纳的绿色调度场景，响应国家绿色调度需求。